\subsection{混合関係数2の$(2,2)$領域には単冪部分がない}
\label{subsec:h3-order2-22-rank2-unipotent}
\begin{lemma}\label{lem:h3-order2-22-rank2-unipotent-rigidity}
$C$をnoetherian quadratic Koszul AS4型14641の整域とし、
$\mu_C$は半単純で固有値$+1,-1$を重複度$(2,2)$で持ち、奇二次関係数は2とする。
$\mu_C$と可換な単冪graded自己同型$H$は恒等である。
ここでは正規三次の存在を仮定しない。
\end{lemma}
\begin{proof}
必要pencil被覆の証明は、半単純中山・Frobenius対合・関係数・整域性だけを使い、
$G$を使っていないので、この$C$にも適用される。
$ruling$型はAS4整域自体から排除済みである。
従ってpure mixed、secant、tangentの三場合を調べればよい。

pure mixedの場合、必要表示の分類部分
（定理\ref{thm:h3-order2-22-pure-exclusion}で正常核を調べる前の部分）により、
二つの係数行列は$p=I,q=\operatorname{diag}(1,z)$、
$z=-1,i,-i$のいずれかである。
$H$は二つの純mixed関係直線を保ち、単冪なので各線を点ごとに固定する。
$H_U=T,H_V=S$と書けば
\[
 TS^t=I,\qquad SqT^t=q.
\]
従って$T^t$は相異なる固有値を持つ$q$と可換し、対角かつ単冪なので$T=I$、$S=I$である。

secantの場合、$E\subset U\otimes V$にはrank-one tensorの直線がちょうど二本ある。
$H$の生成する連結単冪群はこの有限集合を点ごとに固定する。
従って$H_U,H_V$はそれぞれ二軸を保つ対角単冪であり、いずれも恒等である。
この議論はcross blockの階数によらない。

tangentの場合、唯一のrank-one点$xu$を固定するので
\[
 H(x)=x,\quad H(y)=y+\beta x,\quad
 H(u)=u,\quad H(v)=v+\eta u
\]
と書ける。
$E=(xu,xv+yu),F=(ux,uy-vx)$の基底上では
\[
 P=\begin{pmatrix}1&\beta+\eta\\0&1\end{pmatrix},\qquad
 Q=\begin{pmatrix}1&\beta-\eta\\0&1\end{pmatrix},\qquad
 g\longmapsto QgP^{-1}.
\]
可逆cross blockの必要表示では
$g=\left(\begin{smallmatrix}0&-1\\1&0\end{smallmatrix}\right)$であり、
$QgP^{-1}=g$から$\beta=\eta=0$となる。
退化境界のAS4整域として可能な必要表示は、
定理\ref{thm:h3-order2-22-tangent-total-exclusion}の$G$を使わない関係認識部分により
$g=\epsilon I$、正平面は可換、$u^2=\alpha x^2+\delta y^2$、$\delta\ne0$に限る。
この場合はgraph不変性から$\eta=0$、平方関係の不変性から
\[
 \delta(2\beta xy+\beta^2x^2)=0.
\]
$x^2,xy$は可換AS2平面で独立なので$\beta=0$となる。
全ての場合に$H$は全生成元を固定する。
\end{proof}

\begin{theorem}\label{thm:h3-order2-22-rank2-no-Jordan}
noetherian quadratic Koszul AS4型14641の整域$B$の中山半単純部分が
固有値$+1,-1$を重複度$(2,2)$で持ち、その半単純部分による奇二次関係数が2なら、
$\mu_B$自身が半単純である。
従ってこれまでの混合関係数2の被覆・排除は、この重複度の全Jordan型を含む。
\end{theorem}
\begin{proof}
$\mu_B=\mu_s\mu_u$をJordan分解する。
$\mu_u$は$\mu_s$の二つの二次元空間を保つので、
$N=\mu_u|_{B_1}-I$は$N^2=0$を満たす。
$H=I+N/4$と置くと$H^4=\mu_u|_{B_1}$である。
これは$B$のgraded自己同型を定める。
実際、$I+tN$は整数$t$で$\mu_u^t$に等しく二次関係を保つので、
関係保存の多項式恒等式は全$t\in k$で成り立つ。
逆写像は$I-tN$である。

右Zhang twist $C=B^H$を取る。
$H$のKoszul top上の作用も単冪なのでhomological determinantは1であり、
本稿の右twist規約による公式
（\S\ref{sec:as4-h3-reductions}、\cite[Theorem 0.3]{RRZ}）は
\[
 \mu_C=\mu_BH^{-4}=\mu_s
\]
を与える。指数は$-4$である。
$C$もnoetherian quadratic Koszul AS4整域であり、
$H$は$C$の自己同型でもある。
二次関係は$\mathcal R_C=(1\otimes H^{-1})\mathcal R_B$で、
$H$は$\mu_s$と可換なので、奇関係数も2のままである。
$H$は$\mu_C$と可換な単冪自己同型なので、前補題から$H=I$となる。
従って$\mu_u=I$である。
この証明は$G$のtwist下の輸送や、twist後のH3性を仮定していない。
\end{proof}

\begin{corollary}\label{cor:h3-order2-22-rank2-all-Jordan-coverage}
重複度$(2,2)$・奇関係数2の全H3候補は、全Jordan型を含めて、
空である。従って中山半単純部分が重複度$(2,2)$を持つH3包絡では、
奇二次関係の次元は必ず4である。
同じ重複度では、この奇関係数4の全Jordan型が残課題となる。
\end{corollary}
\begin{proof}
前定理で半単純となるので、
定理\ref{thm:h3-order2-22-secant-surface-empty}と既証明の全pencil排除を適用する。
\end{proof}

\texttt{verify\_cubic\_order2\_22\_rank2\_unipotent.py}は
二つのtangent stabilizer、pure mixedの相異なる固有値の中央化条件、
平方関係の不変式、四分根と右twistの$-4$指数を厳密検証する。
必要被覆の$G$非依存性と、有限rank-one点集合への単冪作用の消滅は本文の構造証明である。
追加のas4全被覆定理は輸入していない。
